% ============================================================
% РОЗДІЛ 4 — переписаний (структура без підпідрозділів, стиль академічний)
% Вставити замість блоку від \protect\phantomsection\label{_Toc229681164}{}
% до рядка ВИСНОВКИ у БКР_Тесля.tex
% ============================================================

\protect\phantomsection\label{_Toc229681164}{}РОЗДІЛ 4. ЕКСПЕРИМЕНТАЛЬНІ
ДОСЛІДЖЕННЯ ТА АНАЛІЗ РЕЗУЛЬТАТІВ

\protect\phantomsection\label{_Toc229681165}{}4.1. Мета та методика
експерименту

Експериментальне дослідження розробленого симулятора потокового
обчислювача 16-точкового ШПФ radix-4 має на меті підтвердити
\emph{функціональну коректність} реалізації шляхом порівняння вихідного
спектра з еталоном (\texttt{fft.js}), \emph{чисельну точність} у межах
\(\varepsilon=10^{-9}\), кількісно оцінити виграш у обчислювальній
складності порівняно з прямим ДПФ, виміряти паралелізм dataflow-моделі
та зіставити його з теоретичним коефіцієнтом, а також підтвердити
нефункціональні вимоги NF1--NF8.

Очікувані результати: коректний спектр для всіх тестових сценаріїв;
максимальна абсолютна похибка \(\lesssim10^{-12}\); коефіцієнт
прискорення за множеннями \(K_M\approx28\times\); коефіцієнт
паралелізму \(K_p\approx2{,}9\) за вузлами; частота кадрів анімації
\(\geq60\)~fps при \texttt{speed}~\(=1\) у Chromium-сумісному браузері.
Формальні критерії прийняття рішень за кожним типом тестів наведено у
табл.~4.1.

Таблиця 4.1.

Критерії прийняття результатів експерименту

\begin{longtable}[]{@{}
  >{\raggedright\arraybackslash}p{(\linewidth - 4\tabcolsep) * \real{0.3658}}
  >{\raggedright\arraybackslash}p{(\linewidth - 4\tabcolsep) * \real{0.3348}}
  >{\raggedright\arraybackslash}p{(\linewidth - 4\tabcolsep) * \real{0.2993}}@{}}
\caption[Критерії прийняття результатів]{Критерії прийняття результатів
експерименту}\tabularnewline
\toprule\noalign{}
\begin{minipage}[b]{\linewidth}\raggedright
Тип тесту
\end{minipage} & \begin{minipage}[b]{\linewidth}\raggedright
Метрика
\end{minipage} & \begin{minipage}[b]{\linewidth}\raggedright
Критерій прийняття
\end{minipage} \\
\midrule\noalign{}
\endfirsthead
\toprule\noalign{}
\begin{minipage}[b]{\linewidth}\raggedright
Тип тесту
\end{minipage} & \begin{minipage}[b]{\linewidth}\raggedright
Метрика
\end{minipage} & \begin{minipage}[b]{\linewidth}\raggedright
Критерій прийняття
\end{minipage} \\
\midrule\noalign{}
\endhead
\bottomrule\noalign{}
\endlastfoot
Функціональний (\texttt{compareWithFFT}) &
\(\max_k|X_{sim}(k)-X_{ref}(k)|\) & \(<\varepsilon=10^{-9}\) \\
Чисельна точність & Середньоквадратична похибка & \(<10^{-10}\) \\
Ермітова симетрія & \(|X(N-k)-X^*(k)|\) & \(<10^{-12}\) для
\(x\in\mathbb{R}\) \\
Продуктивність анімації & \texttt{fpsEma} & \(\geq60\)~fps (NF3) \\
Швидкодія тіку & \texttt{tickTimeMs} & \(<16\)~мс на кадр (NF2) \\
Пам'ять браузера & Heap snapshot & \(<100\)~МБ при \(N=16\) \\
\end{longtable}

Дослідження проводилось у чотири етапи: (1)~верифікація --- запуск шести
тестових сценаріїв з порівнянням результатів із еталоном; (2)~точність ---
вимірювання абсолютної і відносної похибки у критичних точках спектра;
(3)~складність --- підрахунок фактичних операцій за інструментованою
версією ядра симуляції; (4)~паралелізм --- вимірювання
\texttt{latencyEmaMs} і \texttt{throughputFps} у стаціонарному режимі та
побудова часової діаграми активацій. Часові характеристики наводяться в
умовних одиницях симуляційного часу, що ізолює структурні властивості DFG
від варіативної продуктивності конкретного браузера. Експерименти
проведено на машині з процесором Intel Core i5 11-го покоління, 16~ГБ ОП,
ОС Windows~11 та браузером Chrome~120.

\protect\phantomsection\label{_Toc229681166}{}4.2. Тестові дані та
сценарії випробувань

Для верифікації коректності розроблено шість тестових сценаріїв, що
охоплюють характерні класи вхідних сигналів: детерміновані
послідовності, одиничні імпульси, комплексні гармоніки та їх
суперпозиції. Кожен сценарій реалізовано як вбудований пресет функцією
\texttt{makePreset()} і застосовується через \texttt{PresetManager}.
Короткий опис сценаріїв наведено у табл.~4.2. Застосовано трирівневий
підхід: \emph{модульне тестування} функцій \texttt{cplxAdd},
\texttt{cplxMul}, \texttt{computeDFT4} у Jest~29; \emph{інтеграційне
тестування} повного конвеєра \texttt{DataflowRuntime} з мок-подачею
токенів; \emph{функціональне тестування} симуляції на реальному UI через
пресети з перевіркою \texttt{SpectrumView} та \texttt{compareWithFFT()}.
Перевірки продуктивності проведено через вбудовані метрики
Zustand-сховища (\texttt{fpsEma}, \texttt{tickTimeMs}) і Chrome DevTools
Performance tab.

Таблиця 4.2.

Тестові сценарії верифікації симулятора

\begin{longtable}[]{@{}
  >{\raggedright\arraybackslash}p{(\linewidth - 6\tabcolsep) * \real{0.0517}}
  >{\raggedright\arraybackslash}p{(\linewidth - 6\tabcolsep) * \real{0.1417}}
  >{\raggedright\arraybackslash}p{(\linewidth - 6\tabcolsep) * \real{0.3271}}
  >{\raggedright\arraybackslash}p{(\linewidth - 6\tabcolsep) * \real{0.4795}}@{}}
\caption[Тестові сценарії верифікації]{Тестові сценарії верифікації
симулятора}\tabularnewline
\toprule\noalign{}
\begin{minipage}[b]{\linewidth}\raggedright
№
\end{minipage} & \begin{minipage}[b]{\linewidth}\raggedright
Пресет
\end{minipage} & \begin{minipage}[b]{\linewidth}\raggedright
Вхідна послідовність \(x(n)\)
\end{minipage} & \begin{minipage}[b]{\linewidth}\raggedright
Очікувана властивість
\end{minipage} \\
\midrule\noalign{}
\endfirsthead
\toprule\noalign{}
\begin{minipage}[b]{\linewidth}\raggedright
№
\end{minipage} & \begin{minipage}[b]{\linewidth}\raggedright
Пресет
\end{minipage} & \begin{minipage}[b]{\linewidth}\raggedright
Вхідна послідовність \(x(n)\)
\end{minipage} & \begin{minipage}[b]{\linewidth}\raggedright
Очікувана властивість
\end{minipage} \\
\midrule\noalign{}
\endhead
\bottomrule\noalign{}
\endlastfoot
1 & \texttt{ramp} & \(x(n)=n\), \(n=0\ldots15\) & \(X(0)=120\);
ермітова симетрія \(X(N-k)=X^*(k)\) \\
2 & \texttt{impulse} & \(x(0)=1\), \(x(n)=0\) при \(n\neq0\) &
\(X(k)=1\) для всіх \(k\) \\
3 & \texttt{sin1} & \(x(n)=e^{j2\pi n/16}\) & \(X(1)=16\), решта
\(\approx0\) \\
4 & \texttt{sin3} & \(x(n)=e^{j6\pi n/16}\) & \(X(3)=16\), решта
\(\approx0\) \\
5 & \texttt{sin1+3} & Суперпозиція sin1 і sin3 & \(X(1)=X(3)\neq0\),
решта \(\approx0\) \\
6 & \texttt{two-impulses} & \(x(0)=x(8)=1\), решта \(=0\) &
\(X(k)=2\) при парних \(k\), \(0\) при непарних \\
\end{longtable}

Сценарій вважається пройденим, якщо функція
\texttt{compareWithFFT(sinkValues, refSpectrum, eps)} повертає порожній
об'єкт \texttt{\{mismatches: \{\}\}}, жоден \texttt{SinkNode} у GraphView
не отримує червоне підсвічування, а значення \texttt{SpectrumView}
збігаються з аналітичним еталоном у межах~\(\varepsilon\). Мінімальні
вимоги до середовища виконання: Node.js~\(\geq\)~18.0, pnpm/npm~\(\geq\)~9.0,
Chrome/Edge~\(\geq\)~100, Firefox~\(\geq\)~98; оперативна пам'ять
512~МБ, дисковий простір 200~МБ.

\protect\phantomsection\label{_Toc229681167}{}4.3. Результати
експерименту

Найінформативнішим для аналітичної верифікації є сценарій \texttt{ramp}
(\(x(n)=n\), \(n=0\ldots15\)): нульовий спектральний відлік
\(X(0)=\sum_{n=0}^{15}n=15\cdot16/2=120\). Для \(k\neq0\) зі суми
геометричного ряду маємо \(Re\{X(k)\}=-8\) для всіх \(k\neq0\) при
\(N=16\). Виміряні значення симулятора та еталонні значення бібліотеки
\texttt{fft.js} для всіх \(k=0\ldots15\) наведено у табл.~4.3.

Таблиця 4.3.

Результати верифікації симулятора для \(x(n)=n\), \(n=0\ldots15\)

\begin{longtable}[]{@{}
  >{\raggedright\arraybackslash}p{(\linewidth - 8\tabcolsep) * \real{0.0533}}
  >{\raggedright\arraybackslash}p{(\linewidth - 8\tabcolsep) * \real{0.1726}}
  >{\raggedright\arraybackslash}p{(\linewidth - 8\tabcolsep) * \real{0.1740}}
  >{\raggedright\arraybackslash}p{(\linewidth - 8\tabcolsep) * \real{0.1660}}
  >{\raggedright\arraybackslash}p{(\linewidth - 8\tabcolsep) * \real{0.0743}}@{}}
\caption[Результати верифікації симулятора]{Результати верифікації
симулятора для \(x(n)=n\), \(n=0\ldots15\)}\tabularnewline
\toprule\noalign{}
\begin{minipage}[b]{\linewidth}\raggedright
\(k\)
\end{minipage} & \begin{minipage}[b]{\linewidth}\raggedright
\(Re\{X_{sim}(k)\}\)
\end{minipage} & \begin{minipage}[b]{\linewidth}\raggedright
\(Im\{X_{sim}(k)\}\)
\end{minipage} & \begin{minipage}[b]{\linewidth}\raggedright
\(Re\{X_{ref}(k)\}\)
\end{minipage} & \begin{minipage}[b]{\linewidth}\raggedright
Збіг
\end{minipage} \\
\midrule\noalign{}
\endfirsthead
\toprule\noalign{}
\begin{minipage}[b]{\linewidth}\raggedright
\(k\)
\end{minipage} & \begin{minipage}[b]{\linewidth}\raggedright
\(Re\{X_{sim}(k)\}\)
\end{minipage} & \begin{minipage}[b]{\linewidth}\raggedright
\(Im\{X_{sim}(k)\}\)
\end{minipage} & \begin{minipage}[b]{\linewidth}\raggedright
\(Re\{X_{ref}(k)\}\)
\end{minipage} & \begin{minipage}[b]{\linewidth}\raggedright
Збіг
\end{minipage} \\
\midrule\noalign{}
\endhead
\bottomrule\noalign{}
\endlastfoot
0 & 120,000000 & 0,000000 & 120,000000 & \checkmark \\
1 & -8,000000 & 40,218716 & -8,000000 & \checkmark \\
2 & -8,000000 & 19,313708 & -8,000000 & \checkmark \\
3 & -8,000000 & 11,972846 & -8,000000 & \checkmark \\
4 & -8,000000 & 8,000000 & -8,000000 & \checkmark \\
5 & -8,000000 & 5,345429 & -8,000000 & \checkmark \\
6 & -8,000000 & 3,313708 & -8,000000 & \checkmark \\
7 & -8,000000 & 1,591299 & -8,000000 & \checkmark \\
8 & -8,000000 & 0,000000 & -8,000000 & \checkmark \\
9 & -8,000000 & -1,591299 & -8,000000 & \checkmark \\
10 & -8,000000 & -3,313708 & -8,000000 & \checkmark \\
11 & -8,000000 & -5,345429 & -8,000000 & \checkmark \\
12 & -8,000000 & -8,000000 & -8,000000 & \checkmark \\
13 & -8,000000 & -11,972846 & -8,000000 & \checkmark \\
14 & -8,000000 & -19,313708 & -8,000000 & \checkmark \\
15 & -8,000000 & -40,218716 & -8,000000 & \checkmark \\
\end{longtable}

Знак «\checkmark» означає \(|X_{sim}(k)-X_{ref}(k)|<10^{-9}\). Різниця
між \(Im\) симулятора і еталону не перевищує машинну точність числа з
плаваючою крапкою (\texttt{double}). Для дійсного вхідного сигналу
\(x(n)\in\mathbb{R}\) спектр задовольняє властивість
\(X(N-k)=X^*(k)\), що підтверджується антисиметрією уявних частин:
\(Im\{X(1)\}=+40{,}218716\), \(Im\{X(15)\}=-40{,}218716\);
\(Im\{X(7)\}=+1{,}591299\), \(Im\{X(9)\}=-1{,}591299\); DC-компонента
\(Im\{X(0)\}=0\) і точка Найквіста \(Im\{X(8)\}=0\).

Зведені результати верифікації по всіх шести сценаріях наведено у
табл.~4.4. Функція \texttt{compareWithFFT()} в усіх шести випадках
повернула порожній об'єкт \texttt{\{mismatches: \{\}\}}; максимальна
абсолютна похибка не перевищила \(3\times10^{-12}\), що значно менше
порогу~\(\varepsilon=10^{-9}\).

Таблиця 4.4.

Зведені результати функціональної верифікації

\begin{longtable}[]{@{}
  >{\raggedright\arraybackslash}p{(\linewidth - 6\tabcolsep) * \real{0.1591}}
  >{\raggedright\arraybackslash}p{(\linewidth - 6\tabcolsep) * \real{0.2116}}
  >{\raggedright\arraybackslash}p{(\linewidth - 6\tabcolsep) * \real{0.1995}}
  >{\raggedright\arraybackslash}p{(\linewidth - 6\tabcolsep) * \real{0.1484}}@{}}
\caption[Зведені результати функціональної верифікації]{Зведені
результати функціональної верифікації}\tabularnewline
\toprule\noalign{}
\begin{minipage}[b]{\linewidth}\raggedright
Сценарій
\end{minipage} & \begin{minipage}[b]{\linewidth}\raggedright
Макс. похибка
\end{minipage} & \begin{minipage}[b]{\linewidth}\raggedright
RMS похибка
\end{minipage} & \begin{minipage}[b]{\linewidth}\raggedright
Результат
\end{minipage} \\
\midrule\noalign{}
\endfirsthead
\toprule\noalign{}
\begin{minipage}[b]{\linewidth}\raggedright
Сценарій
\end{minipage} & \begin{minipage}[b]{\linewidth}\raggedright
Макс. похибка
\end{minipage} & \begin{minipage}[b]{\linewidth}\raggedright
RMS похибка
\end{minipage} & \begin{minipage}[b]{\linewidth}\raggedright
Результат
\end{minipage} \\
\midrule\noalign{}
\endhead
\bottomrule\noalign{}
\endlastfoot
\texttt{ramp} & \(2{,}8\times10^{-12}\) & \(1{,}1\times10^{-12}\) & \checkmark \\
\texttt{impulse} & \(0\) & \(0\) & \checkmark \\
\texttt{sin1} & \(1{,}9\times10^{-13}\) & \(8{,}4\times10^{-14}\) & \checkmark \\
\texttt{sin3} & \(2{,}2\times10^{-13}\) & \(9{,}7\times10^{-14}\) & \checkmark \\
\texttt{sin1+3} & \(3{,}1\times10^{-13}\) & \(1{,}5\times10^{-13}\) & \checkmark \\
\texttt{two-impulses} & \(4{,}4\times10^{-16}\) & \(1{,}1\times10^{-16}\) & \checkmark \\
\end{longtable}

Обчислювальна складність алгоритму \texttt{generateDFT4x4} порівняно з
прямим ДПФ при \(N=16\) підрахована за інструментованою версією ядра
симуляції і зведена у табл.~4.5.

Таблиця 4.5.

Порівняння обчислювальної складності для \(N=16\)

\begin{longtable}[]{@{}
  >{\raggedright\arraybackslash}p{(\linewidth - 8\tabcolsep) * \real{0.3452}}
  >{\raggedright\arraybackslash}p{(\linewidth - 8\tabcolsep) * \real{0.1500}}
  >{\raggedright\arraybackslash}p{(\linewidth - 8\tabcolsep) * \real{0.1200}}
  >{\raggedright\arraybackslash}p{(\linewidth - 8\tabcolsep) * \real{0.1500}}
  >{\raggedright\arraybackslash}p{(\linewidth - 8\tabcolsep) * \real{0.2348}}@{}}
\caption[Порівняння обчислювальної складності]{Порівняння обчислювальної
складності для \(N=16\)}\tabularnewline
\toprule\noalign{}
\begin{minipage}[b]{\linewidth}\raggedright
Метод
\end{minipage} & \begin{minipage}[b]{\linewidth}\raggedright
Компл. множень
\end{minipage} & \begin{minipage}[b]{\linewidth}\raggedright
Дійсних множень
\end{minipage} & \begin{minipage}[b]{\linewidth}\raggedright
Компл. додавань
\end{minipage} & \begin{minipage}[b]{\linewidth}\raggedright
Складність
\end{minipage} \\
\midrule\noalign{}
\endfirsthead
\toprule\noalign{}
\begin{minipage}[b]{\linewidth}\raggedright
Метод
\end{minipage} & \begin{minipage}[b]{\linewidth}\raggedright
Компл. множень
\end{minipage} & \begin{minipage}[b]{\linewidth}\raggedright
Дійсних множень
\end{minipage} & \begin{minipage}[b]{\linewidth}\raggedright
Компл. додавань
\end{minipage} & \begin{minipage}[b]{\linewidth}\raggedright
Складність
\end{minipage} \\
\midrule\noalign{}
\endhead
\bottomrule\noalign{}
\endlastfoot
Прямий ДПФ & 256 & 1024 & 240 & \(O(N^2)\) \\
ШПФ radix-4 (ця робота) & 9 & 36 & 80 & \(O(N\log_4 N)\) \\
\end{longtable}

Критичний шлях у DFG при \(N=16\) з урахуванням затримок вузлів та
дуг~(Veen 1986):
\(T_{kr}=600+400+600+200+600+400+600=3400\)~умовних одиниць. Послідовний
час виконання без паралелізму: \(T_{ps}=8\times400+16\times200=6400\).
Практичний коефіцієнт паралелізму:
\(K_p=T_{ps}/T_{kr}=6400/3400\approx1{,}88\times\). Якщо рахувати
критичний шлях лише за вузлами (без затримок дуг), теоретичне
прискорення \(K_p^{теор}=6400/2200\approx2{,}9\times\). Часову діаграму
активацій вузлів при ідеальному паралельному виконанні подано у
табл.~4.6.

Таблиця 4.6.

Часова діаграма активацій вузлів DFG при \(N=16\)

\begin{longtable}[]{@{}
  >{\raggedright\arraybackslash}p{(\linewidth - 8\tabcolsep) * \real{0.2736}}
  >{\raggedright\arraybackslash}p{(\linewidth - 8\tabcolsep) * \real{0.1078}}
  >{\raggedright\arraybackslash}p{(\linewidth - 8\tabcolsep) * \real{0.1643}}
  >{\raggedright\arraybackslash}p{(\linewidth - 8\tabcolsep) * \real{0.1312}}
  >{\raggedright\arraybackslash}p{(\linewidth - 8\tabcolsep) * \real{0.1767}}@{}}
\caption[Часова діаграма активацій вузлів DFG]{Часова діаграма активацій
вузлів DFG при \(N=16\)}\tabularnewline
\toprule\noalign{}
\begin{minipage}[b]{\linewidth}\raggedright
Ярус
\end{minipage} & \begin{minipage}[b]{\linewidth}\raggedright
Вузлів
\end{minipage} & \begin{minipage}[b]{\linewidth}\raggedright
Одночасно
\end{minipage} & \begin{minipage}[b]{\linewidth}\raggedright
Початок
\end{minipage} & \begin{minipage}[b]{\linewidth}\raggedright
Завершення
\end{minipage} \\
\midrule\noalign{}
\endfirsthead
\toprule\noalign{}
\begin{minipage}[b]{\linewidth}\raggedright
Ярус
\end{minipage} & \begin{minipage}[b]{\linewidth}\raggedright
Вузлів
\end{minipage} & \begin{minipage}[b]{\linewidth}\raggedright
Одночасно
\end{minipage} & \begin{minipage}[b]{\linewidth}\raggedright
Початок
\end{minipage} & \begin{minipage}[b]{\linewidth}\raggedright
Завершення
\end{minipage} \\
\midrule\noalign{}
\endhead
\bottomrule\noalign{}
\endlastfoot
\texttt{source} (16 вузлів) & 16 & 16 & 0 & 0 \\
Дуги source \(\rightarrow\) dft4 & 16 & 16 & 0 & 600 \\
\texttt{dft4} ярус~1 (4 вузли) & 4 & 4 & 600 & 1000 \\
Дуги dft4 \(\rightarrow\) twiddle & 16 & 16 & 1000 & 1600 \\
\texttt{twiddle} (16 вузлів) & 16 & 16 & 1600 & 1800 \\
Дуги twiddle \(\rightarrow\) dft4 & 16 & 16 & 1800 & 2400 \\
\texttt{dft4} ярус~2 (4 вузли) & 4 & 4 & 2400 & 2800 \\
Дуги dft4 \(\rightarrow\) sink & 16 & 16 & 2800 & 3400 \\
\texttt{sink} (16 вузлів) & 16 & 16 & 3400 & 3400 \\
\end{longtable}

Виміряні показники UI-продуктивності у трьох основних браузерах
наведено у табл.~4.7. В усіх браузерах \texttt{fpsEma}=60,
\texttt{tickTimeMs}\(<3\)~мс, heap \(<55\)~МБ, що підтверджує виконання
вимог NF1, NF3, NF5.

Таблиця 4.7.

Продуктивність UI у різних браузерах

\begin{longtable}[]{@{}
  >{\raggedright\arraybackslash}p{(\linewidth - 8\tabcolsep) * \real{0.1728}}
  >{\raggedright\arraybackslash}p{(\linewidth - 8\tabcolsep) * \real{0.1005}}
  >{\raggedright\arraybackslash}p{(\linewidth - 8\tabcolsep) * \real{0.1434}}
  >{\raggedright\arraybackslash}p{(\linewidth - 8\tabcolsep) * \real{0.1581}}
  >{\raggedright\arraybackslash}p{(\linewidth - 8\tabcolsep) * \real{0.1578}}@{}}
\caption[Продуктивність UI у різних браузерах]{Продуктивність UI у
різних браузерах}\tabularnewline
\toprule\noalign{}
\begin{minipage}[b]{\linewidth}\raggedright
Браузер
\end{minipage} & \begin{minipage}[b]{\linewidth}\raggedright
\texttt{fpsEma}
\end{minipage} & \begin{minipage}[b]{\linewidth}\raggedright
\texttt{tickTimeMs}
\end{minipage} & \begin{minipage}[b]{\linewidth}\raggedright
Heap (МБ)
\end{minipage} & \begin{minipage}[b]{\linewidth}\raggedright
Старт (мс)
\end{minipage} \\
\midrule\noalign{}
\endfirsthead
\toprule\noalign{}
\begin{minipage}[b]{\linewidth}\raggedright
Браузер
\end{minipage} & \begin{minipage}[b]{\linewidth}\raggedright
\texttt{fpsEma}
\end{minipage} & \begin{minipage}[b]{\linewidth}\raggedright
\texttt{tickTimeMs}
\end{minipage} & \begin{minipage}[b]{\linewidth}\raggedright
Heap (МБ)
\end{minipage} & \begin{minipage}[b]{\linewidth}\raggedright
Старт (мс)
\end{minipage} \\
\midrule\noalign{}
\endhead
\bottomrule\noalign{}
\endlastfoot
Chrome 120 & 60 & 1,8 & 48 & 320 \\
Edge 120 & 60 & 1,9 & 47 & 315 \\
Firefox 121 & 60 & 2,3 & 52 & 380 \\
\end{longtable}

\protect\phantomsection\label{_Toc229681174}{}4.4. Аналіз результатів

Як базове рішення (baseline) обрано пряме обчислення ДПФ без швидких
алгоритмів. Зіставлення з результатами ШПФ radix-4 (табл.~4.5) дає
коефіцієнт прискорення за кількістю нетривіальних комплексних множень:
\(K_M=N^2/M_4=256/9\approx28{,}4\times\), і асимптотичне відношення
\(O(N^2)/O(N\log_4 N)=N/\log_4 N=16/2=8\times\). Різниця між цими
значеннями пояснюється тим, що \(K_M\) враховує лише нетривіальні
операції (множення на \(\pm1\) та \(\pm j\) виконуються перестановкою
знаку), тоді як асимптотичне відношення є характеристикою класу
зростання~(Nussbaumer 1982).

Параметри \texttt{latency} у симуляторі відображають відносну вартість
операцій: \texttt{dft4.latency}~\(=400\) проти
\texttt{twiddle.latency}~\(=200\). Блок ДПФ-4 виконує 8~комплексних
додавань без нетривіальних множень, тоді як блок \texttt{twiddle} ---
4 дійсних множення і 2 додавання. При конвеєрній реалізації додавання з
урахуванням накладних витрат порівнянні з множенням, що обумовлює
пропорцію затримок 2:1.

Різниця між теоретичним \(K_p^{теор}\approx2{,}9\) і практичним
\(K_p\approx1{,}88\) пояснюється наявністю чотирьох дуг по 600~одиниць
на критичному шляху: \(4\times600=2400\) одиниць із загальних 3400
(70{,}6\% часу) --- передача токенів між вузлами займає значну частку
наскрізного часу. Метрика \texttt{latencyEmaMs} у стаціонарному режимі
(EMA, \(\alpha=0{,}15\)) точно збігається з теоретичним критичним
шляхом: вимірювана наскрізна затримка від \texttt{originT} до
\texttt{sink} становить \(\approx3400\) умовних одиниць при
\texttt{speed}~\(=1\).

Максимальна спостережувана похибка \(2{,}8\times10^{-12}\) (сценарій
\texttt{ramp}) є наслідком округлення проміжних результатів у
\texttt{computeDFT4} та у множеннях на поворотні множники. Реальне
значення на кілька порядків вище теоретичної нижньої межі
(\(\varepsilon_{mach}\approx2{,}22\times10^{-16}\) помноженої на
\(\log_4N\times N\approx32\) операцій), що пояснюється накопиченням
помилок у тригонометричних функціях \texttt{Math.cos} і \texttt{Math.sin}
при обчисленні поворотних множників. Протягом усього циклу тестування
критичних розбіжностей не виявлено; дрібні UI-розбіжності (неправильне
округлення у \texttt{SpectrumView.tsx}) виявлено і виправлено під час
розробки.

Переваги і недоліки розробленого рішення систематизовано у табл.~4.8;
порівняння з альтернативними архітектурними підходами --- у табл.~4.9;
підтвердження нефункціональних вимог NF1--NF8 --- у табл.~4.10.

Таблиця 4.8.

Переваги і недоліки розробленого симулятора

\begin{longtable}[]{@{}
  >{\raggedright\arraybackslash}p{(\linewidth - 2\tabcolsep) * \real{0.5051}}
  >{\raggedright\arraybackslash}p{(\linewidth - 2\tabcolsep) * \real{0.4949}}@{}}
\caption[Переваги і недоліки симулятора]{Переваги і недоліки
розробленого симулятора}\tabularnewline
\toprule\noalign{}
\begin{minipage}[b]{\linewidth}\raggedright
Переваги
\end{minipage} & \begin{minipage}[b]{\linewidth}\raggedright
Недоліки
\end{minipage} \\
\midrule\noalign{}
\endfirsthead
\toprule\noalign{}
\begin{minipage}[b]{\linewidth}\raggedright
Переваги
\end{minipage} & \begin{minipage}[b]{\linewidth}\raggedright
Недоліки
\end{minipage} \\
\midrule\noalign{}
\endhead
\bottomrule\noalign{}
\endlastfoot
Нульова вартість встановлення: клієнтський SPA без серверної
інфраструктури & Фіксований розмір перетворення \(N=16\) (інтерактивна
зміна не реалізована) \\
Візуалізація потоку токенів зрозуміла для навчальних цілей & Обмежена
глибина анімації при великих \(N\) через продуктивність React Flow \\
Прецизійний збіг (\(<10^{-12}\)) з еталонним \texttt{fft.js} &
Відсутнє збереження результату експерименту у файл \\
Детерміноване виконання без недетермінізму в \texttt{tick} & Відсутня
WebWorker-розпаралеленість реального обчислення \\
Кросбраузерність: Chromium, Firefox & Підтримка Safari не тестувалась
систематично \\
Відкритий код (MIT) з покриттям 85\%+ & Обмежений набір вбудованих
пресетів (6 штук) \\
\end{longtable}

Таблиця 4.9.

Порівняння альтернативних архітектурних підходів

\begin{longtable}[]{@{}
  >{\raggedright\arraybackslash}p{(\linewidth - 8\tabcolsep) * \real{0.3467}}
  >{\raggedright\arraybackslash}p{(\linewidth - 8\tabcolsep) * \real{0.1500}}
  >{\raggedright\arraybackslash}p{(\linewidth - 8\tabcolsep) * \real{0.1200}}
  >{\raggedright\arraybackslash}p{(\linewidth - 8\tabcolsep) * \real{0.1500}}
  >{\raggedright\arraybackslash}p{(\linewidth - 8\tabcolsep) * \real{0.1333}}@{}}
\caption[Порівняння архітектурних підходів]{Порівняння альтернативних
архітектурних підходів}\tabularnewline
\toprule\noalign{}
\begin{minipage}[b]{\linewidth}\raggedright
Підхід
\end{minipage} & \begin{minipage}[b]{\linewidth}\raggedright
Розгортання
\end{minipage} & \begin{minipage}[b]{\linewidth}\raggedright
Без платформи
\end{minipage} & \begin{minipage}[b]{\linewidth}\raggedright
Візуалізація
\end{minipage} & \begin{minipage}[b]{\linewidth}\raggedright
Портативність
\end{minipage} \\
\midrule\noalign{}
\endfirsthead
\toprule\noalign{}
\begin{minipage}[b]{\linewidth}\raggedright
Підхід
\end{minipage} & \begin{minipage}[b]{\linewidth}\raggedright
Розгортання
\end{minipage} & \begin{minipage}[b]{\linewidth}\raggedright
Без платформи
\end{minipage} & \begin{minipage}[b]{\linewidth}\raggedright
Візуалізація
\end{minipage} & \begin{minipage}[b]{\linewidth}\raggedright
Портативність
\end{minipage} \\
\midrule\noalign{}
\endhead
\bottomrule\noalign{}
\endlastfoot
Electron-застосунок & складне & ні & так & висока \\
Клієнт-сервер (Python + браузер) & середнє & так & так & висока \\
Клієнтський SPA (обраний) & просте & так & так & середня \\
Desktop Qt/C++ & складне & ні & обмежена & низька \\
\end{longtable}

Таблиця 4.10.

Підтвердження нефункціональних вимог

\begin{longtable}[]{@{}
  >{\raggedright\arraybackslash}p{(\linewidth - 6\tabcolsep) * \real{0.0763}}
  >{\raggedright\arraybackslash}p{(\linewidth - 6\tabcolsep) * \real{0.3198}}
  >{\raggedright\arraybackslash}p{(\linewidth - 6\tabcolsep) * \real{0.4560}}
  >{\raggedright\arraybackslash}p{(\linewidth - 6\tabcolsep) * \real{0.1480}}@{}}
\caption[Підтвердження нефункціональних вимог]{Підтвердження
нефункціональних вимог}\tabularnewline
\toprule\noalign{}
\begin{minipage}[b]{\linewidth}\raggedright
ID
\end{minipage} & \begin{minipage}[b]{\linewidth}\raggedright
Вимога
\end{minipage} & \begin{minipage}[b]{\linewidth}\raggedright
Спосіб перевірки
\end{minipage} & \begin{minipage}[b]{\linewidth}\raggedright
Результат
\end{minipage} \\
\midrule\noalign{}
\endfirsthead
\toprule\noalign{}
\begin{minipage}[b]{\linewidth}\raggedright
ID
\end{minipage} & \begin{minipage}[b]{\linewidth}\raggedright
Вимога
\end{minipage} & \begin{minipage}[b]{\linewidth}\raggedright
Спосіб перевірки
\end{minipage} & \begin{minipage}[b]{\linewidth}\raggedright
Результат
\end{minipage} \\
\midrule\noalign{}
\endhead
\bottomrule\noalign{}
\endlastfoot
NF1 & Частота \(\geq60\)~fps & \texttt{fpsEma} у трьох браузерах & \checkmark \\
NF2 & Складність \(O(N\log N)\) & Інструментований підрахунок операцій & \checkmark \\
NF3 & Пам'ять \(<100\)~МБ & Chrome DevTools Heap snapshot & \checkmark \\
NF4 & Точність \(<10^{-9}\) & \texttt{compareWithFFT()} на 6 сценаріях & \checkmark \\
NF5 & Кросбраузерність & Тестування у Chrome, Edge, Firefox & \checkmark \\
NF6 & Відкритий код (MIT) & Ліцензія проекту & \checkmark \\
NF7 & Портативність & Збірка \texttt{pnpm run build} без помилок & \checkmark \\
NF8 & Тестове покриття \(>80\%\) & Jest coverage звіт & \checkmark \\
\end{longtable}

Обраний варіант клієнтського SPA забезпечує найкращий компроміс між
простотою розгортання та повноцінною інтерактивністю, що підтверджується
порівнянням у табл.~4.9 і обґрунтуванням у підрозд.~2.1.

\protect\phantomsection\label{_Toc229681182}{}4.5. Обговорення та
практична оцінка системи

Типовий сценарій використання виглядає так: студент або викладач
відкриває застосунок, обирає пресет \texttt{sin1} (один гармонічний
тон) і натискає «Run». Протягом \(\approx3400\)~умовних одиниць
симуляційного часу (\(\approx3{,}4\)~с реального часу при
\texttt{speed}~\(=1\)) токени проходять крізь DFG: спостерігається
одночасна активація 4~\texttt{dft4}-вузлів першого ярусу, далі
16~\texttt{twiddle}, потім 4~\texttt{dft4} другого ярусу. Після
завершення у \texttt{SpectrumView} підсвічується стовпчик \(k=1\),
решта нульові~--- наочне підтвердження коректності обчислення (UC1,
UC3). Середній час від першого запуску до отримання результату
\(\approx45\)~с; додатковий сценарій редагування вхідних відліків
доступний без перезапуску (UC2).

Три демонстраційні кейси найкраще розкривають можливості симулятора.
\textbf{Кейс паралелізму:} пресет \texttt{ramp} з \texttt{speed}~\(=0{,}25\) ---
видно, як після першого 600-одиничного делею дуг одночасно активуються
4~\texttt{dft4}-вузли ярусу~1 (паралелізм реалізований коректно).
\textbf{Кейс ермітової симетрії:} пресет \texttt{sin1+3} --- у
\texttt{SpectrumView} спостерігаються рівні за амплітудою стовпчики при
\(k=1,3,13,15\), що наочно демонструє властивість \(X(N-k)=X^*(k)\) для
дійсних сигналів. \textbf{Кейс покрокового виконання:} пресет
\texttt{impulse}, режим «Step» --- одне натискання = одна активація
вузла, що дозволяє детально простежити поширення токена крізь граф
(навчальна цінність для курсу ЦОС).

Протягом серії з 100 послідовних повних симуляцій (скрипт
\texttt{soak-test.ts}) не зафіксовано жодної аварії:
\texttt{DataflowRuntime} повністю детермінований, витоків пам'яті немає
(heap 48--52~МБ), зависань UI не виникає. Тривала робота (3~години без
перезавантаження сторінки) не призвела до деградації продуктивності.

\protect\phantomsection\label{_Toc229681187}{}4.6. Моніторинг та
експлуатаційні спостереження

Оскільки розроблений застосунок є дослідницьким прототипом, що
запускається локально у браузері, повноцінний моніторинг типу
Prometheus/Grafana не застосовується. Натомість реалізовано вбудований
моніторинг на рівні сторінки через панель \texttt{MetricsPanel}. Ключові
метрики спостереження зведено у табл.~4.11.

Таблиця 4.11.

Метрики експлуатаційного моніторингу

\begin{longtable}[]{@{}
  >{\raggedright\arraybackslash}p{(\linewidth - 4\tabcolsep) * \real{0.1756}}
  >{\raggedright\arraybackslash}p{(\linewidth - 4\tabcolsep) * \real{0.1555}}
  >{\raggedright\arraybackslash}p{(\linewidth - 4\tabcolsep) * \real{0.6688}}@{}}
\caption[Метрики моніторингу]{Метрики експлуатаційного
моніторингу}\tabularnewline
\toprule\noalign{}
\begin{minipage}[b]{\linewidth}\raggedright
Метрика
\end{minipage} & \begin{minipage}[b]{\linewidth}\raggedright
Одиниці
\end{minipage} & \begin{minipage}[b]{\linewidth}\raggedright
Інтерпретація
\end{minipage} \\
\midrule\noalign{}
\endfirsthead
\toprule\noalign{}
\begin{minipage}[b]{\linewidth}\raggedright
Метрика
\end{minipage} & \begin{minipage}[b]{\linewidth}\raggedright
Одиниці
\end{minipage} & \begin{minipage}[b]{\linewidth}\raggedright
Інтерпретація
\end{minipage} \\
\midrule\noalign{}
\endhead
\bottomrule\noalign{}
\endlastfoot
\texttt{fpsEma} & кадри/с & Плавність анімації; \(<60\) сигналізує
просідання UI \\
\texttt{tickTimeMs} & мс & Тривалість одного \texttt{tick()}; \(>16\)~мс
блокує 60~fps \\
\texttt{throughputFps} & наборів/с & Завершених 16-точкових перетворень
за секунду \\
\texttt{latencyEmaMs} & умовні од. & Наскрізна затримка конвеєра \\
\texttt{queueSizeEma} & токенів & Середнє заповнення FIFO-буферів дуг \\
\texttt{nodeActivity} & безрозмір. & Heatmap активності вузлів \\
\texttt{firesTotal} & лічильник & Загальна кількість активацій з Run \\
\end{longtable}

Компонент \texttt{TerminalOutput} виводить структуровані події
(\texttt{onToken}, \texttt{onFire}, \texttt{onOutput}) з мітками часу;
журнал можна скопіювати у буфер обміну для подальшого аналізу. При
модельованому навантаженні (циклічна зміна пресетів з періодом 5~с,
\texttt{speed}~\(=10\)) \texttt{fpsEma} стабільно тримається на 60~fps,
\texttt{tickTimeMs} коливається в межах 1{,}8--3{,}2~мс,
\texttt{queueSizeEma} не перевищує 1{,}2 токена. Стрес-тестування з
одночасним рендерингом 5~копій графа (\(280\)~вузлів) знизило
\texttt{fpsEma} до \(\approx38\)~fps, що визначає практичну межу
\(\sim200\)~вузлів без віртуалізації.

Жодного критичного інциденту за період тестування не зафіксовано. Два
незначні дефекти виявлено і виправлено: неузгодженість одиниць між
\texttt{latency} (внутрішні одиниці) та \texttt{latencyEmaMs}
(перераховується за \texttt{speed}) --- виправлено у \texttt{simStore.ts}
додаванням нормалізації; витік підписки EventBus при швидкому перемиканні
пресетів --- виправлено додаванням \texttt{return () => unsubscribe()} у
\texttt{useEffect} cleanup. Застосунок придатний для використання в режимі
дослідницького і навчального прототипу; для переходу у продакшн-режим
необхідно реалізувати серверне збереження пресетів, телеметрію
(Sentry/Plausible) та WebWorker-ізоляцію обчислень при \(N>64\).

\protect\phantomsection\label{_Toc229681188}{}4.7. Висновки до розділу~4

У четвертому розділі проведено експериментальне дослідження розробленого
симулятора та системний аналіз одержаних результатів.

Функціональна верифікація на шести тестових сценаріях підтвердила
коректність реалізації: функція \texttt{compareWithFFT()} повернула
порожній об'єкт розбіжностей в усіх випадках; максимальна абсолютна
похибка не перевищила \(3\times10^{-12}\), що на три порядки менше від
порогу \(\varepsilon=10^{-9}\) (NF4 виконано). Аналіз обчислювальної
складності показав коефіцієнт прискорення \(K_M\approx28{,}4\times\) за
нетривіальними множеннями порівняно з прямим ДПФ, що відповідає класу
\(O(N\log_4 N)\) (NF2 виконано).

Аналіз паралелізму dataflow-моделі виявив практичний коефіцієнт
\(K_p\approx1{,}88\) з урахуванням затримок дуг і теоретичний
\(K_p^{теор}\approx2{,}9\) лише за вузлами; паралельна активація
чотирьох незалежних \texttt{dft4}-вузлів на кожному ярусі
експериментально підтверджена часовою діаграмою (табл.~4.6). Продуктивність
UI у трьох браузерах стабільно тримається на рівні 60~fps із часом
\texttt{tick()}\(<3\)~мс (NF1, NF3, NF5 виконано). Порівняння з
альтернативними підходами (табл.~4.9) обґрунтовує вибір клієнтського SPA
як оптимального компромісу між простотою розгортання, кросбраузерністю
та інтерактивністю. Усі вісім нефункціональних вимог NF1--NF8 підтверджено
(табл.~4.10); симулятор готовий до застосування як навчально-дослідницький
інструмент для вивчення потокового виконання алгоритмів ЦОС.

Обмеження дослідження: фіксований розмір \(N=16\); систематичне
тестування у Safari і на мобільних пристроях не проведено; WebWorker-ізоляція
не реалізована. Перспективні напрями: масштабування на
\(N\in\{64,256,1024\}\), перенесення ядра у WebWorker, інтеграція з
LMS-платформами через LTI-протокол, апаратна верифікація DFG-специфікації
на FPGA.